%%
%% Test: RTL/Bidi Support — v1.1
%% Stage: xepersian refactor (stages 1-4)
%% 
%% Purpose: Verify that xepersian is loaded correctly:
%%   - xepersian is the LAST package in matinbook.cls
%%   - Persian fonts work (XB Niloofar by default)
%%   - RTL and LTR environments work
%%   - Persian numbering works
%%   - Latin environment works
%%   - All previous modules still work
%% 

\documentclass{matinbook}

\title{تست پشتیبانی RTL — نسخه ۱.۱}
\author{تیم MatinBook}
\date{۱۴۰۵}

\begin{document}

\maketitle

\chapter{بررسی پشتیبانی RTL/Bidi}
\label{chap:test-rtl}

این سند، بارگذاری درست \texttt{xepersian} و پشتیبانی RTL/Bidi
را بررسی می‌کند.

%----------------------------------------
\section{متن فارسی}
\label{sec:persian-text}

این یک متن فارسی است. اعداد فارسی: ۱۲۳۴۵۶۷۸۹۰

شماره‌گذاری فصل‌ها و بخش‌ها باید به فارسی باشد.

%----------------------------------------
\section{محیط لاتین}
\label{sec:latin-env}

\begin{latin}
This is a Latin paragraph inside an RTL document.
Numbers: 1234567890
Special characters: ! @ \# \$ \% \& * ( )
\end{latin}

%----------------------------------------
\section{متن دوزبانه}
\label{sec:bilingual}

در اینجا یک کلمه‌ی لاتین \lr{Hello World} در متن فارسی داریم.

و در اینجا یک کلمه‌ی فارسی در متن لاتین:

\begin{latin}
The Persian word \rl{سلام} means ``hello''.
\end{latin}

%----------------------------------------
\section{ریاضیات}
\label{sec:math}

معادله‌ی زیر در محیط RTL:

\begin{equation}
    \label{eq:rtl-test}
    f(x) = \sum_{i=1}^{n} x_i^2
\end{equation}

نمادها: $\mathbb{N} \subset \mathbb{Z} \subset \mathbb{Q} \subset \mathbb{R} \subset \mathbb{C}$

%----------------------------------------
\section{قضیه}
\label{sec:theorem}

\begin{theorem}[قضیه‌ی نمونه]
    برای هر عدد حقیقی $x$، داریم $x^2 \geq 0$.
\end{theorem}

\begin{proof}
    چون مربع هر عدد حقیقی نامنفی است، نتیجه حاصل می‌شود.
\end{proof}

%----------------------------------------
\section{کد}
\label{sec:code}

\begin{latin}
\begin{minted}{python}
def greet(name):
    """Return a greeting."""
    return f"Hello, {name}!"
\end{minted}
\end{latin}

%----------------------------------------
\section{جدول}
\label{sec:table}

\begin{table}[h]
\centering
\caption{جدول نمونه}
\label{tab:rtl-test}
\begin{tabular}{L{3cm} C{2cm} R{3cm}}
\toprule
\textbf{ستون ۱} & \textbf{ستون ۲} & \textbf{ستون ۳} \\
\midrule
مقدار ۱ & ۱۰ & مقدار ۳ \\
مقدار ۲ & ۲۰ & مقدار ۴ \\
\bottomrule
\end{tabular}
\end{table}

%----------------------------------------
\section{ارجاعات}
\label{sec:references}

ارجاع به معادله: \cref{eq:rtl-test}

ارجاع به جدول: \cref{tab:rtl-test}

ارجاع به فصل: \cref{chap:test-rtl}

%----------------------------------------
\section{نتیجه‌گیری}
\label{sec:conclusion}

اگر این سند بدون خطا کامپایل شود:

\begin{itemize}
    \item ✅ \texttt{xepersian} آخرین پکیج در \texttt{matinbook.cls} است
    \item ✅ فونت XB Niloofar درست بارگذاری شده
    \item ✅ محیط \texttt{latin} کار می‌کند
    \item ✅ \texttt{\textbackslash lr} و \texttt{\textbackslash rl} کار می‌کنند
    \item ✅ شماره‌گذاری فارسی کار می‌کند
    \item ✅ ریاضیات، قضیه، کد، جدول، ارجاعات کار می‌کنند
\end{itemize}

\textbf{وضعیت تست:} قبول

\end{document}
